\documentclass{article}
\usepackage{amsmath,amssymb,amsthm,url}
\newtheorem{theorem}{Theorem}
\newtheorem{proposition}{Proposition}
\newtheorem{lemma}{Lemma}
\title{Eager activation and common cuts: the finite-causal-past boundary}
\author{}\date{17 November 2026}
\begin{document}\maketitle
\paragraph{Audit and declared change of interface.}
This is a structural scheduling probe, not an extension of the old deterministic same-clock native crossing. The observation is an occurrence-labelled execution diagram together with a chosen sequential realization; comparisons may change the realization order. Executed event identities, raw records, paid-event order, original birth sponsors and whole-law commitments remain, whereas timestamps and the order of mutually exchangeable silent instructions need not. Same-clock full-prefix equality, unchanged early-port availability, unrestricted controller transport and arbitrary current-world seeds are explicitly outside the conclusion. The theorem below distinguishes a finite scheduling repair from an infinite barrier; stochastic enlargement alone cannot repair the latter.

The pinned baseline is repository revision \texttt{88e05c7d526474bedaef1cab19d53cc7ab127acf}, especially \emph{Beyond Guard Projective Iteration and Crossing Obstruction}, \emph{World Rebased Multi Active Open Assembly}, and the three nonblocking recurrence manuscripts. Those sources permit unrelated ready work without proving schedule interchange. They retain actual world arrows and explicitly reject automatic seed naturality. The new claim therefore starts with an independently restricted exchangeable presentation; it does not silently add interchange to the baseline.

\section{The graph question, before probabilities or effects}
Fix one fully specified occurrence diagram $E$. Its elements are actual instruction occurrences, not types or output values. It is finite or countably infinite, and has a strict prerequisite relation $\prec$ closed transitively. Its edges include producer/read dependencies, the prescribed order of all paid instructions, and any fixed internal whole-block order. A finite event has finitely many direct input ports, but its transitive predecessor set may nevertheless be infinite. No finite-branching or finite-indegree assumption is substituted for finite transitive past.

A sequential complete realization is an injective enumeration of every event, finite when $E$ is finite and of order type $\omega$ otherwise, in which each prerequisite precedes its consumer. Completeness is a fairness condition on this fixed actual diagram: every event is eventually executed. It is not fairness of every alternative command or of unchosen hypothetical events. All raw outcomes, including zero-probability outcomes, have their own diagrams.

A proposed common-cut discipline adds directed edges $B$ to $\prec$. For example, to require a continuation $y$ to wait for the return of a whole upstream forest, add $x\to y$ for every event $x$ in that actual upstream forest, including recursive unfold instructions. Let $\prec_B$ be the transitive closure. A cycle is already an obstruction. An event with infinitely many $\prec_B$ predecessors is a different obstruction even when the graph has no cycle.

\begin{theorem}[Exact pathwise barrier-realizability criterion]\label{thm:criterion}
The enriched prerequisite diagram admits a complete sequential realization if and only if it is acyclic and every event has finite transitive past. This criterion applies to finite and countable diagrams, including diagrams with infinitely many roots and no uniform bound on finite causal depths.
\end{theorem}
\begin{proof}
In an enumeration indexed by natural numbers, every event at position $n$ has at most $n$ predecessors, so both conditions are necessary. For sufficiency fix any enumeration $e_0,e_1,\ldots$ of the event set. Repeatedly choose the least-index unexecuted event whose predecessors have all been executed. Such an event exists whenever events remain: choose one unexecuted event, and a minimal element of its finite unexecuted downset. Every event is eventually chosen. Prove this by induction on the size of its finite downset. Its finitely many predecessors are eventually chosen by the induction hypothesis. Once it is eligible, only finitely many lower-index events can be chosen ahead of it; no event is chosen twice. Thus it is chosen after finitely many additional choices. No bound on the sizes of other downsets is used.
\end{proof}

The theorem is an existence and explicit countable scheduling construction on a supplied diagram. It is not an effective procedure deciding an infinite graph's finite-past property. In a dynamically generated implementation the scheduler additionally needs an available finite certificate that a proposed gate is satisfied. A positive completed-forest receipt supplies such a certificate. The theorem alone does not give an oracle predicting future return.

\begin{proposition}[Infinite barrier obstruction independent of randomness]\label{prop:infinite}
If an actual locally ready continuation $y$ is behind a common cut containing infinitely many executed upstream events, no complete sequential realization of that same occurrence diagram can enforce that cut. This remains true with arbitrary random tapes, noncausal choice of tapes, and any law on the tapes, whenever each sample path still consists of discrete finite instructions indexed by natural numbers and preserves the event occurrences and prerequisite meaning.
\end{proposition}
\begin{proof}
On each sample path preserving $y$, every upstream event must appear before its finite position. There are infinitely many such events, contradicting finiteness of that position's past. Randomization cannot change this pathwise contradiction. Suppressing $y$, changing its prerequisite set, or introducing a transfinite activation are changes of interface, not solutions within the stated class.
\end{proof}

This strengthens the old missing-fair-bit witness in a different direction. That witness excludes deterministic exact raw pushforward even without a causality requirement. Proposition~\ref{prop:infinite} permits additional random resources but requires preservation of the actual event and cut structure; it excludes a full event-preserving schedule repair even for deterministic $y$. The two obstructions are independent: one concerns available random law, the other finite execution positions. Neither proves that all useful observation-indexed or stochastic comparisons fail.

\section{When the graph criterion transports the actual machine}
The graph criterion does not itself justify exchanging a paid update with initialization. A reusable sufficient primitive subclass is the following. Paid occurrences retain their original total order and actual world arrows. Only a silent instruction $s$ and a paid instruction $p$ with no read, creation, completion or resource dependency between them may be exchanged. The paid instruction cannot select an actor newly born at $s$, inspect its fresh certificate or use its fresh raw value. Its symbolic command and paired commitment are the same on both routes. The full per-step label contains a different currently revealed reserve on the two routes and is NOT asserted literally identical; its label is reconstructed from the same original commitment and the appropriate revealed prefix. The silent readiness predicate factors through its immutable birth/binder fields and the retained procedural counter, with every required read already present. Its passive response to the paid update preserves those fields and counter; readiness therefore remains true after the update. The paid eligibility and selected-row domain factor through the world and the pre-existing selected actor fields, which the silent instruction does not modify. Its raw-total compatibility predicates require actual passive responses on every possible newly installed silent state/frame; these responses are supplied on the entire two-step raw carrier, including zeros. Thus no world-dependent gate, missing new-state response, global quiescence flag, or fresh-certificate predicate is silently classified as independent. These local domain conditions, together with the raw-domain snapshot equations, derive legality of both orders. The silent instruction is an original snapshot/reveal instruction whose immutable raw block, binder and commitment do not depend on whether $p$ has occurred. Any newly installed state is transported by its supplied passive map along the actual intervening world path. Already installed actors are never reset.

Two silent instructions are exchangeable only when neither creates or reads the other's fields, they update disjoint code/registry entries, each readiness and row-domain predicate factors through fields left unchanged by the other, and either their owners differ or the intrinsic factor ordering explicitly permits their labelled exchange. Their endpoint equation follows from disjoint record union and the same labelled tensor/cone calculation. Paid instructions are totally ordered and are never exchanged with one another.

For law comparisons one retains, at the common pre-swap boundary, a SINGLE specified whole joint of the matched raw events and its original test projections, or a specified causally executable resource realization of that joint. The law is not replaced by independent primitive marginals. Each ordering executes conditionals of that same commitment, including its separately specified order-indexed zero tables. This is a comparison of complete committed plans, not automatic permission to insert an early commitment instruction into an existing implementation. The actual per-step incoming law is the posterior on the revealed prefix; using the same old incoming marginal after reserve-dependent selection would violate calibration. At positive fibres the finite chain rule supplies the correct posterior and service. A claimed implementation must expose the paired commitment and the required reads as an authorized interface, and retain the actual residual labels; this draft does not supply such a realization for arbitrary controllers. General original local primitive families are not retested against newly enlarged execution histories. A source-full-binder law may push into a target-local-binder family by the original convex-mixture argument, but the converse and equality of all-law envelopes are not asserted.

For material, $s$ inserts its actual factor into an already existing historical owner slot; $p$ creates a later disjoint paid head scope. The primitive factor and new head are held fixed. The update to the older owner expression commutes with tensoring the new scope factor, using the labelled tensor symmetry and the source's actual insertion cone. Thus both endpoint expressions have the same original owner slots, not the arrival-ordered tensor. For several factors of one owner, retain their fixed intrinsic consumption order as graph edges; arbitrary factor permutations are not included. Full content grades and unused shapes are carried by identity/readdressing. The swap changes exposure time, so equality of the full timestamped birth registry is deliberately not claimed.

\begin{lemma}[Derived local snapshot exchange]\label{lem:exchange}
For the subclass just stated, the two routes $s;p$ and $p;s$ have a canonical bijection of matched complete raw two-step records and the same endpoint installed/pending registry after readdressing occurrence identities. They carry the same specified joint law and actual endpoint genealogical material by the supplied passive maps and insertion cones. The two intermediate observation states and early exposure times can differ.
\end{lemma}
\begin{proof}
On both routes the paid instruction has the same old actor tuple, raw paid witness and actual update $u$. Its active posts are identical because it cannot inspect or select the silent-born fields. Every pre-existing unselected actor and frame receives its same passive response. If the silent instruction installs a snapshot state $z$ before $u$, that state becomes $R_u z$; if it installs afterward, the snapshot declaration installs precisely $R_u z$. All immutable binder, commitment, sponsor and consumed-block fields agree. The same reasoning applies to a factor-free call-unfold whose code state is preserved by passive transport. No active/passive interchange for arbitrary active commands was used. Sampling in either order by the conditional tables of the same complete joint gives that joint by the finite chain rule. All complete raw zero assignments remain retained records; the order-indexed zero-prefix execution tables are explicit fields and are not claimed determined by the joint. Thus equality of the complete joint does not identify the full intermediate zero-table or label interfaces. The two material maps agree at every old owner projection and at the new head projection; tensor coherence and uniqueness of the owner cone prove the endpoint equation.
\end{proof}

This is not available for a current-world initializer caching the world: initialize before an increment and passively preserve the cache gives $0$, while initialize afterward gives $1$. That baseline witness violates the snapshot premise exactly. A world-receptive pending frame alone is insufficient. A controller reading the newly produced bit before selecting $p$ also violates the independence premise, even if the two marginal raw laws happen to agree.

\begin{theorem}[Finite exchange and its unbounded limit]\label{thm:exchange}
Call a finite downset replay-admissible when its retained raw records reconstruct a legal finite prefix, its world-sensitive records include their actual required paid prefix and passive paths, and its discarded factors have the baseline augmentation restrictions. In the fixed snapshot/read-independent presentation above this is checked by the stated local gate/domain rules and the retained prerequisite edges, not by a claimed reverse world arrow. On a fixed diagram whose incomparable occurrences have the applicable primitive exchanges, any two finite linearizations of the same finite down-closed event set are connected by adjacent exchanges. Their endpoint native fields and the specified whole joint are carried by the composite local maps. If two complete countable realizations exist, each finite replay-admissible down-closed observation is reconstructible from both, after possibly different finite numbers of instructions; the resulting event-indexed finite native cones agree. This yields equality only at this declared event-indexed observation, not a same-clock full-prefix map.
\end{theorem}
\begin{proof}
For finite linearizations, move the first event of the second ordering leftward in the first ordering. Every crossed event is incomparable with it, otherwise one of the two orders violates a prerequisite. Apply its adjacent exchange, then induct on the remaining finite set. All maps are determined by the same occurrence-labelled raw assignment, actual passive path and owner-cone projections; hence different exchange words induce the same endpoint field map. For two countable realizations and a finite replay-admissible down-closed set $F$, completeness gives a finite time by which all events of $F$ have occurred in each. In each larger finite prefix, use the finite exchange result to move all events of $F$ to the front; down-closure ensures that no prerequisite is crossed. Read the resulting $F$ prefix from the retained historical raw records. Its material restriction is the actual augmentation of exactly the omitted factors, using the original owner cones. This is a reconstruction of an event-indexed history observation, NOT an inverse world arrow or a map from the later current actor state to an earlier one. The finite result identifies the $F$ component, and restriction equations on inclusions of finite downsets follow by uniqueness of reconstruction and the augmentation equations. No infinite product of swap maps or last step at time $\omega$ is taken.
\end{proof}

The replay-admissibility qualification in the theorem is material: omitted paid events cannot be forgotten while retaining a state at their later world as if the world had not advanced. One safe choice is to include the entire initial segment of paid events required by every retained state, together with all their actual prerequisites, and to retain passive path fields. If these closures are infinite, that requested observation is not finite and the theorem makes no claim about it. Likewise an event-indexed cone does not identify the original controller tables or preserve first-done times.

\section{Compositional meaning and boundary examples}
For finitely many component diagrams, gluing is union of their occurrence and dependency fields with matching shared names, not independent resampling. A locally finite-past discipline is not compositional merely because each component passes the test. Two components can contribute edges $x\to y$ and $y\to x$ and create a cycle. More subtly, one component may have a chain $x_0\prec x_1\prec\cdots$ and another introduce $x_n\to y$ for all $n$; the union is acyclic but gives $y$ infinite past. These are distinct failure witnesses. The admitted global gluing domain is therefore the unions whose actual generated relation is acyclic with finite past, together with the primitive exchange/data conditions. Associativity of union preserves the result for an already admissible full union; it does not turn independently admissible pieces into an admissible union.

The positive example has two recursive upstream branches, each of which may have arbitrarily long but finite actual initialization. A local continuation after branch one may run early. On every raw diagram in which both forests are finite, adding a common return cut gives finite past and Theorem~\ref{thm:criterion} gives a barrier realization; snapshot/read-independence then gives the finite event-aligned comparison. The common cut must be fixed without permitting a newly created continuation to influence the upstream forests; otherwise changing the schedule changes the actual diagram and this example does not meet the premise. No uniform bound on initialization length is needed. If one raw branch is an infinite unfold chain, the same continuation acquires infinite past and the comparison loses that event. An almost-sure termination statement cannot remove this all-raw boundary: a zero-probability infinite branch still fails the criterion if raw preservation is required.

\paragraph{What this settles and what remains.}
The result identifies an exact order-theoretic obstruction and an independently supplied primitive regime in which schedule reordering really carries native endpoints and committed laws. It supplies a discriminating test for proposed eager/common-cut bridges and explains why finite exchange proofs do not establish an unbounded barrier repair. It does not establish the old deterministic full-prefix crossing, universal stochastic native cells, arbitrary current-world exchange, controller-view equivalence, early-port contextual equivalence, first-done clock preservation or a largest feedback doctrine. A useful full eager extension must either stay inside the finite-past/exchangeable interface or explicitly retain the changed gates, observation and random resources. The event diagram is supplied at each full raw assignment; deriving a common occurrence diagram from arbitrary history-dependent controllers remains a separate obligation. In particular, the positive scheduling construction is not by itself an authorized-view controller. Those obligations remain visible rather than being absorbed into a statement that eager execution is merely a faster schedule.
\end{document}
